\chapter{Introduction}
\label{ch:introduction}

\section{Motivation}
With the widespread use of the internet, online shopping has become very common. However, size charts, model photos and static clothes images often lead to uncertainty, making users have difficulty intuitively imagining how these online clothes they buy will look on them. Furthermore, the existing virtual try-on applications are few, and even if they try to solve this problem, they cannot reflect the actual wardrobe situation of users due to their single functionality. At the same time, many clothing recommendation systems only provide general suggestions, ignoring users' styles, occasions and personal clothing collections (See Chapter \ref{chapter: Existing Market Applications} for details).

The starting point of this project lies in the gap between user needs and existing solutions. By integrating multiple functions such as virtual try-on, personalized clothing recommendations and users' wardrobes, we aim to provide users with a more practical dressing decision-making system.

\section{Aims and Objectives}
This project aims to develop an AI-powered wardrobe assistant that enables users to manage clothing items, explore virtual outfit combinations, and receive personalized styling recommendations. By taking advantage of LLM and model-based image editing, the system intends to provide an accessible decision-support tool that improves  the efficiency of daily dressing and enhances users' fashion experience.

The research objectives are as follows:
\begin{itemize}
    \item \textbf{Wardrobe management system:} Design and implement a cross-platform wardrobe management module that supports clothes upload, classification, metadata editing, and user profile binding.
    \item \textbf{Digital wardrobe dataset:} Construct a structured wardrobe dataset containing key clothes attributes (e.g. category, color, material, season) and possible descriptive semantics for personalization tasks.
    \item \textbf{Context-aware recommendations:} Integrate LLM to generate outfit suggestions that take into account user preferences, historical choices, weather conditions, and social occasions.
    \item \textbf{Virtual try-on pipeline:} Implement an interactive virtual try-on feature that enables users to generate outfit previews and compare different clothing combinations.
    \item \textbf{Interpretability and trend awareness:} Provide transparent recommendations and trend-aware suggestions, helping users understand styling logic and make informed, controllable wardrobe decisions.
\end{itemize}